% Refer back to paper syntax and definitons - Don't repeat what has already been written instead use what has been defined as basis and show important differences in how the formalization differes from the paper syntax.

% This chapter is supposed to be where background binding representation and the $\pi$LL syntax comes together to creates a mechanization of the pre established theory. Should probably include:
  % How Processes are defined in Lean and how LN has influences the constructors
  % How Types are defined in Lean and how De Bruijn influced the constructors
  % How Processes and Types are used together to define the Typing system, as well as what influence LN and De Bruijn had on the constructors in the form of including co-finite quantification.
  % Define the operational semantics for the Typing sytem and how it can be mapped to an operational semantic for Processes and Types, resepctively.
  % Mention the elimination of Alpha equivalence, (Would be good to have a proof that it has been reduced to structural equivalence)

% When talking about the formlization, maybe show a Lean version of a mapping to the reports running example, and comment on how, due to it being a strict application, there are sometimes empty hyperenv or env artifacts which roam around the derivation. Contextualize this in process congruence and the monoidal structure of typing environments, and mention that this property has yet to be added to the formalization.
  % Maybe something along the lines of:
  % Notice that in the intermediate transition, the strict application of the \textsc{Mix} rule yields the typing context $\emptyset \mid y:1$. By maintaining the structural footprint of the terminated left-hand process ($\vdash \textbf{0} :: \emptyset$), the exact correspondence with the parallel transition rule remains visually explicit. Following the final $y[]$ transition, the process fully terminates as $\textbf{0} \mid \textbf{0}$. At this stage, we can apply process congruence ($P \mid \textbf{0} \equiv P$) alongside the monoidal identity laws of the hyperenvironment ($\emptyset \mid \emptyset = \emptyset$) to cleanly resolve the derivation to its final state: $\vdash \textbf{0} :: \emptyset$


% ---------- SPLIT ---------- %

% Write about Core system, syntax overview, process structure, judgements, operational semantics etc. with focus on the multiplicative component.

% 3. πLL System
    % 3.1 Syntax (full system)
    % 3.2 Types (full system)
    % 3.3 Typing (full system)
    % 3.4 Operational Semantics (multiplicative fragment)

% Generally
    % Present syntax for processes and types
    % Show some Typing rules
    % For operational semantics, show TypingStep 
        % Define process and types semantics 
        % Make this very breif


% \section{$\pi$MLL: Paper Level Definitions and Syntax}
% ------- WIP
% \section{Operational Semantics}

% \paragraph{Fragment decomposition.}
% The system \pill can be naturally decomposed into two complementary fragments. The first is the multiplicative fragment, \pimll, which includes constructs such as tensor, par, and their associated units. This fragment captures the core structure of interaction, describing how processes communicate and synchronize through linear channels. The second is the non-multiplicative fragment, comprising additives, exponentials, and quantifiers. These constructs enrich the language with additional expressive power, enabling features such as internal and external choice, controlled duplication and disposal of resources, and polymorphic session behavior.

% In this development, the syntactic and typing components are defined for the full \pill system, encompassing both fragments. However, the operational semantics and associated metatheory are established for the multiplicative fragment, which suffices to capture the essential communication behavior underlying the system. Extending the operational treatment to the non-multiplicative fragment introduces additional considerations, such as choice resolution and resource management, and is left for future work.

% --- MERGE THE TWO THINGS ABOVE --- %

% Aim to bridge theory-code gap and explain design choices
  % 1. Design goals
  % 2. LN for processes (how does this eliminates AlphaEq)
  % 3. De Bruijn indices for types (needed type variables in first iteration to be locally closed, so instead of assuming this as hypothesis, implemented De Bruijn indices and the notion of a type being locally closed based on the current binder level for a typing derivation then 'n' in front of judgements)
  % 4, Co-finite quantification 